\documentclass[a4paper,twoside]{article}

\usepackage{morewrites}

\usepackage[svgnames,dvipsnames,x11names]{xcolor}
\usepackage{graphicx}
\usepackage{texsmith-lists}
\usepackage[english]{babel}
\usepackage[colorlinks=true, linkcolor=NavyBlue, urlcolor=NavyBlue, citecolor=NavyBlue]{hyperref}
\usepackage[a4paper]{geometry}
\usepackage{ts-fonts}
\usepackage{float}
\usepackage{graphicx}
\usepackage{amsmath}
\usepackage{seqsplit}
\usepackage{ts-code}

\usepackage{lastpage}
\usepackage{refcount}
\makeatletter
\def\ps@plain{
\let\@oddhead\@empty
\let\@evenhead\@empty
\def\@oddfoot{
\hfil
\ifnum\getpagerefnumber{LastPage}>1\relax
\thepage
\fi
\hfil}
\let\@evenfoot\@oddfoot
}
\makeatother

\title{Snippet Blocks}
\author{}\date{}
\begin{document}
\maketitle
TeXSmith renders fenced blocks with the \tscodeinline{.snippet} class into PDF/PNG pairs and injects them into the page with a download link. Snippets now accept a concise YAML payload instead of \tscodeinline{data-\allowbreak{}*} attributes, and the PNG preview is left unframed so you can wrap it with the \tscodeinline{ts-\allowbreak{}frame} fragment when needed.
\section{YAML-driven snippets}
Use a YAML fence to point TeXSmith at your sources, working directory, and extra template options. The width of the rendered preview comes from the fence attribute.
\begin{tscode}[lang=md, engine=pygments]
\begin{Verbatim}[commandchars=\\\{\},breaklines, breakanywhere, commandchars=\\\{\}]
\PY{l+s+sb}{```}\PY{l+s+sb}{yaml}\PY{+w}{ }\PYZob{}.snippet\PYZcb{}
\PY{n+nt}{layout}\PY{p}{:}\PY{+w}{ }\PY{l+lScalar+lScalarPlain}{2x2}
\PY{n+nt}{cwd}\PY{p}{:}\PY{+w}{ }\PY{l+lScalar+lScalarPlain}{../../examples/paper}
\PY{n+nt}{sources}\PY{p}{:}
\PY{+w}{  }\PY{p+pIndicator}{\PYZhy{}}\PY{+w}{ }\PY{l+lScalar+lScalarPlain}{cheese.md}
\PY{+w}{  }\PY{p+pIndicator}{\PYZhy{}}\PY{+w}{ }\PY{l+lScalar+lScalarPlain}{cheese.bib}
\PY{n+nt}{template}\PY{p}{:}\PY{+w}{ }\PY{l+lScalar+lScalarPlain}{article}
\PY{n+nt}{width}\PY{p}{:}\PY{+w}{ }\PY{l+lScalar+lScalarPlain}{70\PYZpc{}}
\PY{n+nt}{fragments}\PY{p}{:}
\PY{+w}{  }\PY{l+lScalar+lScalarPlain}{ts\PYZhy{}frame}
\PY{n+nt}{press}\PY{p}{:}
\PY{+w}{  }\PY{n+nt}{frame}\PY{p}{:}\PY{+w}{ }\PY{l+lScalar+lScalarPlain}{true}
\PY{l+s+sb}{```}
\end{Verbatim}
\end{tscode}
\begin{itemize}
\item \tscodeinline{cwd} is the directory where TeXSmith runs, resolving relative sources.
\item \tscodeinline{sources} mirrors the CLI arguments: Markdown inputs and auxiliary files like \tscodeinline{.bib}.
\item \tscodeinline{layout} arranges multiple PDF pages on the preview grid.
\item \tscodeinline{press} merges into the template context (fragments, format, etc.).
\end{itemize}
\section{Inline Markdown with front matter}
You can also render inline Markdown. Front matter drives the template choice and fragment selection while the body becomes the snippet content.
\begin{tscode}[lang=md, engine=pygments]
\begin{Verbatim}[commandchars=\\\{\},breaklines, breakanywhere, commandchars=\\\{\}]
\PY{g+gu}{```md \PYZob{}.snippet caption=\PYZdq{}Inline snippet\PYZdq{} width=\PYZdq{}65\PYZpc{}\PYZdq{}\PYZcb{}}
\PY{g+gu}{\PYZhy{}\PYZhy{}\PYZhy{}}
template: snippet
fragments:
  ts\PYZhy{}frame:
press:
\PY{g+gu}{  frame: true}
\PY{g+gu}{\PYZhy{}\PYZhy{}\PYZhy{}}
\PY{g+gh}{\PYZsh{} Section}

Some content...
```
\end{Verbatim}
\end{tscode}
\begin{figure}[H]
\centering
\includegraphics[width=0.65\linewidth]{snippet-<HASH>.pdf}
\caption[Inline snippet]{Inline snippet}
\end{figure}
\section{Reusing a configuration file}
When the same snippet settings are shared across pages, point the fence to a YAML config and keep the block body empty or minimal.
\begin{tscode}[lang=md, engine=pygments]
\begin{Verbatim}[commandchars=\\\{\},breaklines, breakanywhere, commandchars=\\\{\}]
\PY{l+s+sb}{```}\PY{l+s+sb}{md}\PY{+w}{ }\PYZob{}.snippet config=\PYZdq{}snippet\PYZhy{}configs/letter.yml\PYZdq{} caption=\PYZdq{}Config\PYZhy{}driven snippet\PYZdq{} width=\PYZdq{}70\PYZpc{}\PYZdq{}\PYZcb{}
\PY{l+s+sb}{```}
\end{Verbatim}
\end{tscode}
The configuration lives alongside this page at \tscodeinline{docs/examples/snippet-\allowbreak{}configs/letter.yml}.
\end{document}
